% V2第20章原书思考题20.1--20.10，PDF物理页204，独立候选；图20.20--20.21待精确重绘。
\section*{第20章\quad 电磁感应：思考题}
\paragraph{思考题 20.1}
一只灵敏电流计的线圈在永磁体磁场中偏转后，若将线圈两端短接会很快停转。解释这是什么电磁现象。

\paragraph{思考题 20.2}
高频炉的感应线圈里通高频交流电为何能加热放在其中的金属块？热量的能量源是什么？

\paragraph{思考题 20.3}
变压器铁芯为什么做成相互绝缘的薄叠片，而不采用一整块厚铁芯？叠片平面的方向与铁芯中的交变磁通有何关系？

\paragraph{思考题 20.4}
尺寸完全相同的铜环和铝环置于相同变化磁通中，两环感生电动势与感生电流是否相等？明确电阻与自感的影响。

\paragraph{思考题 20.5}
电子感应加速器里的电子从何处获得动能？用法拉第定律和电场对电荷做功解释。

\paragraph{思考题 20.6}
三个彼此靠近的线圈中心在同一条直线上，怎样安放中间线圈使它与两边的互感尽可能接近零？

\paragraph{思考题 20.7}
大小不同的两个闭合线圈相互靠近，为使互感尽量大，应该将小线圈套在大线圈内还是把它们面对面放置？讨论条件。

\paragraph{思考题 20.8}
电路有强电流，断开开关时出现跨越触点的火花。用自感现象解释电压突升。

\paragraph{思考题 20.9}
解释图20.20超导磁悬浮体在外磁场中受到稳定悬浮力的原因，指出理想抗磁性与磁通钉扎的区别。

\paragraph{思考题 20.10}
金属探测器能够识别地下金属物体，为何其探头内的交流线圈可通过涡电流效应产生可测的返回信号？恒定直流是否同样有效？
